\section{DIMACS and Variable Mapping Output} \label{sec:dimacs-and-variable-mapping-output}
The system exports two important files for verification and debugging.

\subsection{DIMACS CNF} \label{subsec:dimacs-cnf}

The generated SAT instance is exported in DIMACS CNF format.
The file contains:

\begin{itemize}
    \item comment lines beginning with \texttt{c},
    \item a problem line of the form \texttt{p cnf <variables> <clauses>},
    \item one line per clause, ending in \texttt{0}.
\end{itemize}

\noindent An example clause line may look like:
\begin{lstlisting}[label={lst:clause-example}]
-10 -274 26 0
\end{lstlisting}

\noindent This allows the generated formula to be inspected manually or passed to an external SAT solver.

\subsection{Variable Mapping} \label{subsec:variable-mapping}

The variable mapping translates integer SAT variable identifiers into readable names.
For example:
\begin{lstlisting}[label={lst:mapping-example}]
{
  "274": "ascends(Alice, 0)",
  "275": "descends(Alice, 0)",
  "710": "carries(Alice, Sofa, 0)"
}
\end{lstlisting}

\noindent This mapping is essential for interpreting a satisfying truth assignment.